Read in Gauge Theory Gravity, the metric used is the derived metric $g(b,c) = \hb^{-1}(b)\cdot\hb^{-1}(c) = b\cdot c - 2H(b\cdot k)(c\cdot k)$ of the Kerr-Schild gauge field of Fact \claimref{interior.imported-kerr-schild-form-kerr-gauge} (using $Q(k) = 0$), written in components in the signature $(-,+,+,+)$, and induced metrics are Gram matrices of physical tangent vectors. All entries are rational in $(M, a, r, u)$.

\emph{(i)} The coordinate $u$ does not occur in $k$ and $\eta$ is diagonal, so $g_{uu} = \Sigma/(1-u^2)$ and $g_{tu} = g_{\phi u} = 0$; the Gram determinant on $r = \mathrm{const}$ is $g_{uu}$ times the determinant $D$ of the $(t,\phi)$ block. Replacing $\partial_t$ by $\partial_t + \Omega\partial_\phi$ adds a multiple of one row and column to another and leaves $D$ unchanged. The block is $\eta_b + 2H\,k_b\otimes k_b$ with $\eta_b = \mathrm{diag}(-1, P(1-u^2))$ and $k_b = (-1, a(1-u^2))$, so by the rank-one determinant rule $D = \det\eta_b\,(1 + 2H\,k_b\cdot\eta_b^{-1}k_b) = -P(1-u^2)\bigl(1 + 2H(-1 + a^2(1-u^2)/P)\bigr) = -(1-u^2)\bigl(P - 2H(P - a^2 + a^2u^2)\bigr) = -(1-u^2)(P - 2H\Sigma) = -(1-u^2)(P - 2Mr) = -(1-u^2)\Delta$. Hence the Gram determinant is $-\Delta\Sigma$; the volume element per unit Killing parameter is $\sqrt{|\Delta|\Sigma}\,du\,d\phi$, and integrating over $\phi \in [0,2\pi)$ and $u \in [-1,1]$ gives $V'(r)$. The sign of $-\Delta\Sigma$ gives the causal character (in this signature a positive Gram determinant is a spacelike three-surface). At $a = 0$, $\Sigma = r^2$, $I = 2r$, $\Delta = r^2(1 - 2M/r)$. The check confirms the determinant at rational points with three angular velocities each.

\emph{(ii)} The inverse metric is $\eta^{-1} - 2H\,k^\sharp\otimes k^\sharp$ with $k^\sharp$ the $\eta$-raised $k$, whose $u$ component vanishes and whose $r$ component is $-1$; so $g^{ru} = 0$ and $g^{rr} = P/\Sigma - 2H = \Delta/\Sigma$ (the stored relation $\Sigma g^{rr} = \Delta$). The gradient of $r$ therefore has no $u$ component, its lines run at fixed $u$, and the proper length along them is $\int dr/\sqrt{|g^{rr}|} = \int\sqrt{\Sigma/|\Delta|}\,dr$. Near $r_+$ write $\delta = |r - r_+|$: $|\Delta| = \delta\,(2s \mp \delta)$ and $\Sigma = \Sigma_+(u)(1 + O(\delta))$ with $\Sigma_+ = r_+^2 + a^2u^2 > 0$, so $\ell_+ = \sqrt{\Sigma_+/2s}\int_0^{\delta}x^{-1/2}(1 + O(x))\,dx = 2\sqrt{\Sigma_+\delta/2s}\,(1 + O(\delta))$ and $w = \sqrt{2s\,\delta\,\Sigma_+}\,(1 + O(\delta))$. Hence $w/(P_+\ell_+) \to 2s/(2P_+) = (r_+ - r_-)/(2P_+) = \kappa_+$, in which $\Sigma_+(u)$ has cancelled: the limit is the same at every latitude. The same computation at $r_-$, where $|\Delta| = \delta(2s \mp \delta)$ again and $\Sigma_- = r_-^2 + a^2u^2 > 0$ for $a > 0$, gives $\kappa_- = (r_+ - r_-)/(2P_-)$. That $P_\pm\,du\,d\phi$ is the area element of the cross-section: on $r = r_\pm$ at fixed $t$ the Gram determinant of $(\partial_u, \partial_\phi)$ is $g_{uu}g_{\phi\phi}$ with $g_{\phi\phi} = P(1-u^2) + 2Ha^2(1-u^2)^2$, so $g_{uu}g_{\phi\phi} = P\Sigma + 2Mr\,a^2(1-u^2)$, and at $\Delta = 0$, where $2Mr = P$, this is $P(\Sigma + a^2 - a^2u^2) = P^2$. The check verifies the constant term of the polynomial $|\Delta|\Sigma/\delta$ in $\delta$ against $\kappa_\pm^2P_\pm^2\cdot 4\Sigma_\pm/2s$ at four latitudes for five holes, and $P_\pm = 2Mr_\pm$.

\emph{(iii)} $\mathcal{S}_{\rho_0}$ is $r = r_+ + \delta(u)$ (or $r_+ - \delta(u)$) with $\ell_+ = \rho_0$, so $\delta(u) = \rho_0^2\,2s/(4\Sigma_+(u))\,(1 + O(\delta))$. Its tangent vectors are $T = \partial_t + \Omega\partial_\phi$, $U = \partial_u + \delta'(u)\partial_r$ and $\partial_\phi$, with $\delta' = O(\delta)$ since $\delta$ is a smooth multiple of $1/\Sigma_+(u)$. Against the Gram matrix of (i) only the entries of $U$ change: $g(T,U) = \delta'g(T,\partial_r)$, $g(U,\partial_\phi) = \delta'g_{r\phi}$, $g(U,U) = g_{uu} + \delta'^2g_{rr}$, with the components $g_{r\mu}$ bounded near the locus in this gauge. Expanding along $U$, the determinant is $g(U,U)\,D$ minus a quadratic form in $(g(T,U), g(U,\partial_\phi))$ with bounded coefficients, that is $-\Delta\Sigma\,(1 + O(\delta^2)) + O(\delta^2)$; since $\Delta\Sigma$ is of order $\delta$, the determinant is $-\Delta\Sigma\,(1 + O(\delta))$. Hence $V'_{\rho_0} = \int\!\!\int w\,(1 + O(\delta))\,du\,d\phi = \kappa_+\rho_0\int\!\!\int P_+\,du\,d\phi\,(1 + O(\delta)) = \kappa_+\rho_0A_+(1 + O(\delta))$ by (ii). All lengths scale with $M$ at fixed $a/M$, and $\delta/M = O(\rho_0^2/M^2)$ at fixed $a/M < 1$, which gives the limit. The count is $\nu V'_{\rho_0}\beta_+$ by the rule of Proposition \claimref{singularity.stretched-singularity-count} with the epoch of Theorem \claimref{scrambling.epoch-modular-period-two-unknowns-become}; $\kappa_+$ cancels, and $2\pi\nu\rho_0 = \tfrac14$ at $\nu\rho_0 = 1/8\pi$. The entropy $A_+/4$ compared with is that of Fact \claimref{synthesis.imported-kerr-surface-gravity-entropy}.

\emph{(iv)} On $r = \mathrm{const}$ the Gram determinant of (i), taken with $\Omega = \Omega_H$, is $g_{uu}\bigl(g(K,K)\,g_{\phi\phi} - g(K,\partial_\phi)^2\bigr)$ in components. On the locus $g(K, \partial_\phi) = 0$ (there $D = 0$ and $g(K,K) = 0$, so $g(K,\partial_\phi)^2 = 0$), hence $g(K,\partial_\phi) = O(\delta)$ nearby and $|Q(K)| = w^2/(g_{uu}g_{\phi\phi})\,(1 + O(\delta)) = (w/P_+)^2(1 + O(\delta))$. So (ii) says $\sqrt{|Q(K)|} = \kappa_+\ell_+(1 + o(1))$ with $\kappa_+$ independent of $u$: the surface gravity of Fact \claimref{synthesis.imported-kerr-surface-gravity-entropy} read as the growth of the Killing norm per unit proper distance from the horizon (the classical equivalent form of $(K\cdot D)K = \kappa_+K$ for a Killing horizon, used here as a reading and verified directly by (ii) for this field), together with its constancy on the horizon. The count of (iii) uses nothing else about the hole, so no value of $(M, a)$ with $a < M$ could give a ratio to $A_+/4$ other than $8\pi\nu\rho_0$ at leading order.

\emph{(v)} With $\beta_+ = 2\pi/\kappa_+ = 4\pi Mr_+/s$ and $A_+/4 = 2\pi Mr_+$: rate $8\pi\nu M$ gives $N = 32\pi^2\nu M^2r_+/s = 2\pi M^2r_+/s$, ratio $M/s$; rate $\nu\beta_+$ gives $N = \beta_+^2/16\pi = \pi M^2r_+^2/s^2$, ratio $Mr_+/2s^2$; rate $2\nu\kappa_+A_+$ gives $4\pi\nu A_+ = A_+/4$. With $e = a/M$, $\sigma = s/M$, $1 - \sigma = e^2/(1+\sigma)$: $M/s - 1 = e^2/(\sigma(1+\sigma))$ and $Mr_+/2s^2 - 1 = (1-\sigma)(1+2\sigma)/2\sigma^2 = e^2(1+2\sigma)/(2\sigma^2(1+\sigma))$, with coefficients $\tfrac12$ and $\tfrac34$ at $\sigma = 1$; both ratios diverge as $\sigma \to 0$. At $a = 0$ all three are $1$, which is the stored Schwarzschild count. For the last display: $T_HS = (\kappa_+/2\pi)(\pi P_+) = s/2$, and $2\Omega_HJ = 2a^2M/(2Mr_+) = a^2/r_+ = r_-$, so $M - 2\Omega_HJ = s = 2T_HS$ (the Smarr relation, verified here directly), and $2\nu\kappa_+A_+ = 8\pi\nu s$.

\emph{(vi)} $\kappa_\pm A_\pm = 4\pi P_\pm(r_+ - r_-)/(2P_\pm) = 4\pi s$. For a surface $r = \mathrm{const}$, $N = \nu V'\beta_+ = \tfrac{1}{16\pi}\cdot2\pi\sqrt{|\Delta|}\,I\cdot4\pi Mr_+/s$, and $N = 2\pi Mr_+$ exactly when $\sqrt{|\Delta|}\,I = 4s$. At $a = 0$: $\sqrt{r|2M - r|}\cdot2r = 4M$, i.e.\ $r^3|2M - r| = 4M^2$. On $(0, 2M)$ the function $r^3(2M - r)$ rises to its maximum $27M^4/16$ at $r = 3M/2$ and then falls to zero, so it takes the value $4M^2$ exactly twice when $27M^4/16 > 4M^2$, i.e.\ $M^2 > 64/27$; on $(2M, \infty)$, $r^3(r - 2M)$ increases from $0$ without bound and takes it once. At $r = r_c$, $r_c^3 = 2M$: $r_c^3(2M - r_c) - 4M^2 = -r_c^4 < 0$, and at $r_c(1 + x)$, $x = r_c^{-2}$, it is $r_c^6[(1+x)^3(1 - x - x^2) - 1] = r_c^6\,x\,(2 - x - 5x^2 - 4x^3 - x^4)$, which is positive for $x \le \tfrac14$; so the first root lies in $(r_c, r_c + 1/r_c)$ for $r_c \ge 2$. Near the horizon, $r = 2M \mp d$ gives $(2M \mp d)^3d = 4M^2$, $d = \tfrac{1}{2M}(1 + O(M^{-2}))$, and by (ii) with $\Sigma_+ = 4M^2$, $2s = 2M$ the proper distance is $2\sqrt{4M^2d/2M} = 2\sqrt{2Md} = 2$ at leading order. The leading rates agree: $\kappa\rho_0A = \tfrac{1}{4M}\cdot2\cdot16\pi M^2 = 8\pi M$. The check brackets the three roots exactly at $r_c = 4$ and $r_c = 10$.

\emph{(vii)} The first sentence is (v) with $s \to 0$ and $r_+ \to M$, $A_+ = 8\pi Mr_+ \to 8\pi M^2$. At $a = M$, $\sqrt{\Sigma/|\Delta|} = \sqrt{\Sigma}/|r - M|$ with $\Sigma \ge M^2u^2 + r^2 > 0$ bounded below near $r = M$, so $\int_M^r$ diverges logarithmically; and $\kappa_+ = s/(2Mr_+) = 0$.
